% Owners: Hiền (models and tests), Tân (evaluation and comparison).
% Numbers come from output/inference/ (scripts 03, 04 and 05).
\newcommand{\inferencefigure}[4]{%
  \begin{figure}[H]
    \centering
    \includegraphics[width=0.96\linewidth,height=#4\textheight,keepaspectratio]{figures/inference/#1}
    \caption{#2}
    \label{#3}
  \end{figure}%
}
\section{Analysis Sample and Test Design}
Chapter~5 uses the 2,423 deduplicated records from Chapter~3. A seeded
80/20 split divides them into 1,938 training records (309 ads) and 485
held-out test records (77 ads). Records with an identical feature vector
form one group, and each group goes entirely to one side, so the test set
contains no copy of a training record. All hypothesis tests and model
coefficients use the training records only; the test set is used once, at
the end, to measure how well the models work on new records.

The population behind these tests is the kind of web images that the
original collection sampled. The results describe associations, not causes.

\section{Do Ads and Non-Ads Differ in Width?}
Chapter~4 showed that ads are much wider than non-ads in the sample. We test
whether this difference holds in the population with a two-sample Welch
$t$-test, which does not assume equal variances:
\[
  t=\frac{\bar{x}_{\mathrm{ad}}-\bar{x}_{\mathrm{non}}}
       {\sqrt{s^2_{\mathrm{ad}}/n_{\mathrm{ad}}+s^2_{\mathrm{non}}/n_{\mathrm{non}}}},
  \qquad H_0:\ \mu_{\mathrm{ad}}=\mu_{\mathrm{non}},\quad
  H_1:\ \mu_{\mathrm{ad}}\neq\mu_{\mathrm{non}}.
\]
The test uses the training records with an observed width
(Table~\ref{tab:infer-width}). Width is right-skewed, but each class has
hundreds of observations, so by the central limit theorem the sample means
are close to normal and the test is valid.

\begin{table}[htbp]
\centering
\caption{Observed width by class in the training sample.}
\label{tab:infer-width}
\begin{tabular}{lrrrr}
\toprule
Class & $n$ & Mean (px) & SD (px) & Median (px)\\
\midrule
Ad & 265 & 325.6 & 162.0 & 455\\
Non-Ad & 1,269 & 124.2 & 82.1 & 101\\
\bottomrule
\end{tabular}
\end{table}

The mean difference is 201.5~px, with a 95\% confidence interval of
[181.4, 221.6]~px ($t=19.73$, $\mathrm{df}\approx 292.9$, $p<0.001$). We
reject $H_0$: in the population, ads are wider than non-ads on average, by
roughly 180 to 220 pixels. Section~\ref{sec:infer-or} checks whether width
still matters after the other features are taken into account.

\section{Is Missingness Associated with Advertisement Status?}
Chapter~4 also showed that missing size information is more common among
non-ads. The null hypothesis is that \texttt{missing\_size} (any original
feature missing) and ad status are independent. We use Pearson's chi-square
test without continuity correction on the training table
(Table~\ref{tab:infer-missing}); the smallest expected count is 66.65, well
above five.

\begin{table}[htbp]
\centering
\caption{Missingness and class in the training sample.}
\label{tab:infer-missing}
\begin{tabular}{lrrr}
\toprule
Any feature missing & Ad & Non-Ad & Total\\
\midrule
No & 263 & 1,257 & 1,520\\
Yes & 46 & 372 & 418\\
\bottomrule
\end{tabular}
\end{table}

The ad rate is 11.0\% when a feature is missing and 17.3\% otherwise.
The test gives $\chi^2(1)=9.70$ and $p=0.0018$, so at the 5\% level we
reject independence. The association is weak, however: Cram\'er's
$V=0.071$.

\section{Which Features Are Associated with Ad Odds?}
\label{sec:infer-or}
We fit a logistic regression of ad status on four predictors: $\log_2$
Height, $\log_2$ Width, Local and \texttt{missing\_size}. With base-2 logarithms, the odds
ratio (OR) of a dimension is the change in odds when that dimension doubles
and the other predictors stay fixed. Missing dimensions are filled with
training medians and a missing Local with the training mode, and the flag
\texttt{missing\_size} records that a value was filled. Aspect ratio is
left out because it is computed from width and height. We call this the
\emph{main-effects model}; Section~\ref{sec:infer-compare} also uses it as
the one reference model for all comparisons.

\begin{table}[htbp]
\centering
\caption{Adjusted odds ratios from the main-effects logistic model (training data).}
\label{tab:infer-or}
\begin{tabular}{lrrr}
\toprule
Predictor & OR & 95\% Wald interval & $p$\\
\midrule
Height, doubled & 0.95 & [0.82, 1.10] & 0.497\\
Width, doubled & 5.16 & [4.25, 6.25] & $<0.001$\\
Local = 1 & 1.55 & [1.10, 2.17] & 0.011\\
Any feature missing & 1.45 & [0.99, 2.13] & 0.055\\
\bottomrule
\end{tabular}
\end{table}

Width has by far the strongest association: each doubling of width
multiplies the odds of being an ad by about 5.2 (95\% CI 4.3 to 6.3).
Local also has a positive association (OR 1.55). For height, the interval
contains 1, so there is no evidence of an association once width is in the
model; the missingness flag is also not significant at the 5\% level. Its OR
points in the opposite direction from the raw rates in
Table~\ref{tab:infer-missing}. This can happen after adjustment, because the
model compares a missing record with complete records of the same
filled-in size, not with all complete records.

\section{Do Term Indicators Add Information?}
The 1,554 binary term indicators record whether a phrase appears in the
text around an image. Their names come from the file \texttt{ad.names} in
the UCI distribution of the dataset (DOI 10.24432/C5V011); the Kaggle CSV
keeps only the column numbers. A name has the form
\texttt{field*phrase}: the field shows where the phrase occurs
(\texttt{url} and \texttt{ancurl} are URL fields, where \texttt{ancurl} is
the URL that the image links to; \texttt{alt} is the image's alternative
text), and a plus sign joins adjacent words.

From the training records, we select the six terms with the largest
absolute difference between their prevalence among ads and among non-ads,
considering only terms present in at least 20 training records. Because the
data chose these terms, the tests below describe strong associations in
this dataset; they do not confirm a hypothesis fixed in advance.

\subsection{Marginal association}
For each eligible term we test independence between its 0/1 value and the
class with Pearson's chi-square test. Of the 213 terms present in at least
20 training records, 116 have all expected counts of at least five and are
tested; the Benjamini--Hochberg procedure adjusts their $p$ values for
multiple testing. All six selected terms have adjusted $p<0.001$
(Table~\ref{tab:infer-term-chi}). Each of them is present in 27 to 45\%
of training ads but in at most 1.2\% of non-ads.

\begin{table}[htbp]
\centering
\caption{The six selected terms on training data: number of the 309 ads and 1,629 non-ads that contain each term. Each chi-square test has one degree of freedom.}
\label{tab:infer-term-chi}
\begin{tabular}{llrrrr}
\toprule
Column & Term & Ads & Non-ads & $\chi^2$ & BH $p$\\
\midrule
\texttt{term\_1243} & \texttt{ancurl*com} & 139 & 10 & 720.5 & $<0.001$\\
\texttt{term\_0351} & \texttt{url*ads} & 97 & 4 & 510.0 & $<0.001$\\
\texttt{term\_1399} & \texttt{ancurl*click} & 95 & 4 & 498.4 & $<0.001$\\
\texttt{term\_1483} & \texttt{alt*click} & 92 & 19 & 393.7 & $<0.001$\\
\texttt{term\_1435} & \texttt{alt*here} & 85 & 19 & 354.9 & $<0.001$\\
\texttt{term\_1455} & \texttt{alt*click+here} & 84 & 14 & 374.9 & $<0.001$\\
\bottomrule
\end{tabular}
\end{table}

The terms read like advertising: the image file sits in a path containing
``ads'', the image links to an address containing ``click'', and the alt
text says ``click here''.

\subsection{Association after adjustment}
We add the six terms to the main-effects model. Together they reduce the
training deviance by 523.8 on 6 degrees of freedom (likelihood-ratio test,
$p<0.001$), so the terms carry information beyond size, Local and
missingness. Table~\ref{tab:infer-term-or} shows each term's OR with the
other predictors held fixed.

\begin{table}[htbp]
\centering
\caption{Term odds ratios adjusted for size, Local, missingness and the other five terms (training data).}
\label{tab:infer-term-or}
\begin{tabular}{lrrr}
\toprule
Term & OR & 95\% Wald interval & $p$\\
\midrule
\texttt{ancurl*com} & 50.6 & [21.3, 120.3] & $<0.001$\\
\texttt{url*ads} & 116.9 & [33.2, 412.1] & $<0.001$\\
\texttt{ancurl*click} & 39.1 & [10.6, 143.8] & $<0.001$\\
\texttt{alt*click} & 6.9 & [1.4, 33.7] & 0.017\\
\texttt{alt*here} & 0.91 & [0.04, 19.3] & 0.949\\
\texttt{alt*click+here} & 1.23 & [0.03, 43.7] & 0.909\\
\bottomrule
\end{tabular}
\end{table}

The first four terms keep strong positive associations. The last two have
very wide intervals because the three \texttt{alt} terms carry almost the
same information: of the 117 training records that contain any of them, 98
contain all three, and \texttt{alt*click+here} never occurs without the
other two. Once \texttt{alt*click} is in the model, the other two add almost
nothing that the data can separate, so their own ORs cannot be estimated
precisely (collinearity). The very large ORs also reflect that some
term combinations occur almost only among ads; R reported fitted
probabilities numerically equal to 0 or 1 for this model, so these
estimates are imprecise and should be read as ``very strong'' rather than
as exact values.

\section{Comparing the Models on New Records}
\label{sec:infer-compare}
The main-effects model is the reference. We compare it with two
extensions, each adding one kind of information:
\begin{itemize}
\item \textbf{Format + interaction} replaces the straight-line size effects
  with three-degree-of-freedom natural splines, adds a wide-banner indicator
  (aspect ratio $\geq 6$) and a Height $\times$ Width interaction. It asks
  whether a curved size effect helps.
\item \textbf{Main effects + six terms} adds the six term indicators. It asks
  whether the words around an image help.
\end{itemize}

All three models use the same five grouped training folds, and missing
values are filled inside each fold. For the term model, the six terms are
also selected again inside each fold (in folds 3--5,
\texttt{ancurl*redirect} replaced \texttt{alt*click+here}). Each model's
decision threshold is set on its out-of-fold training predictions: the
threshold with the highest ad recall among those with precision of at least
0.70. The threshold is then frozen and applied once to the test set.

Average precision (AP) summarizes how well a model ranks ads above
non-ads; a model that ranks at random has an expected AP equal to the share
of ads, 0.159 in the test set. Precision and recall describe the decisions
at the frozen threshold. Because the test set has only 77 ads, each test
precision and recall is given with a 95\% Wilson confidence interval.

\begin{table}[htbp]
\centering
\small
\setlength{\tabcolsep}{4pt}
\caption{Model results on out-of-fold training predictions (309 ads) and the held-out test set (77 ads), with 95\% Wilson intervals.}
\label{tab:infer-models}
\begin{tabular}{lrrccrr}
\toprule
Model & AP & ROC-AUC & Precision [95\% CI] & Recall [95\% CI] & FP & FN\\
\midrule
\multicolumn{7}{l}{\emph{Training, out-of-fold predictions (309 ads)}}\\
Main effects & 0.574 & 0.792 & 0.74 [0.67, 0.79] & 0.50 [0.44, 0.55] & 55 & 156\\
Format + interaction & 0.679 & 0.833 & 0.75 [0.69, 0.80] & 0.56 [0.50, 0.61] & 57 & 136\\
Main effects + six terms & 0.843 & 0.920 & 0.70 [0.66, 0.75] & 0.85 [0.81, 0.89] & 111 & 46\\
\midrule
\multicolumn{7}{l}{\emph{Held-out test set (77 ads)}}\\
Main effects & 0.594 & 0.762 & 0.66 [0.53, 0.76] & 0.55 [0.43, 0.65] & 22 & 35\\
Format + interaction & 0.684 & 0.827 & 0.68 [0.56, 0.78] & 0.58 [0.47, 0.69] & 21 & 32\\
Main effects + six terms & 0.795 & 0.895 & 0.61 [0.51, 0.69] & 0.84 [0.75, 0.91] & 42 & 12\\
\bottomrule
\end{tabular}
\end{table}

\inferencefigure{precision_recall_curves.png}{Precision--recall curves of the three models. The dotted line marks the training precision target of 0.70.}{fig:infer-pr}{0.34}

\inferencefigure{test_confusion_matrices.png}{Held-out confusion matrices at the frozen thresholds, with 95\% Wilson intervals for precision and recall.}{fig:infer-confusion}{0.30}

Both extensions are judged against the main-effects model on the same 485
test records. To test whether one model catches more ads, we use an exact
McNemar test: it counts only the ads that one model catches and the other
misses, and tests whether these two counts are equal. The same test
compares false alarms on non-ads. For the difference in AP, a paired
bootstrap (2,000 resamples of the test set, ads and non-ads resampled
separately) gives a 95\% confidence interval.

\begin{table}[htbp]
\centering
\small
\setlength{\tabcolsep}{4pt}
\caption{Paired comparisons with the main-effects model on the test set. ``Main / new'' counts the records flagged by only the main-effects model and by only the new model; $p$ is from the exact McNemar test.}
\label{tab:infer-paired}
\begin{tabular}{lccccc}
\toprule
& \multicolumn{2}{c}{Ads caught (77)} & \multicolumn{2}{c}{Non-ads flagged (408)} & AP gain\\
\cmidrule(lr){2-3}\cmidrule(lr){4-5}
New model & Main / new & $p$ & Main / new & $p$ & [95\% CI]\\
\midrule
Format + interaction & 0 / 3 & 0.25 & 9 / 8 & 1.00 & 0.090 [0.039, 0.145]\\
Main effects + six terms & 0 / 23 & $<0.001$ & 0 / 20 & $<0.001$ & 0.200 [0.092, 0.318]\\
\bottomrule
\end{tabular}
\end{table}

\paragraph{A curved size effect ranks better but decides the same.}
The format + interaction model raises test AP from 0.594 to 0.684, and the
interval of the gain, [0.039, 0.145], excludes zero; the same gain appears
in the training out-of-fold predictions. At the frozen thresholds, however,
it catches only 3 more ads than the main-effects model and flags 1 fewer
non-ad. These differences are well within chance (McNemar $p=0.25$ and
$p=1.00$). The curved size effect improves the ranking of ads, but at this
threshold its decisions are not distinguishable from the simpler model.

\paragraph{The words around an image add the most.}
Adding the six terms raises test AP from 0.594 to 0.795, a gain of 0.200
(95\% CI 0.092 to 0.318). The term model catches every ad the main-effects
model catches plus 23 more: 65 of 77 against 42 of 77 (McNemar
$p<0.001$). Its recall interval, [0.75, 0.91], does not overlap that of the
main-effects model, [0.43, 0.65]. The price is 20 extra false alarms
(42 against 22 non-ads flagged, $p<0.001$), while precision changes from
0.66 to 0.61.

\paragraph{The precision target on new records.}
The thresholds aim for precision of at least 0.70 on training data. For the
main-effects and format + interaction models, the test precision intervals,
[0.53, 0.76] and [0.56, 0.78], contain 0.70, so the test set is too small to
show a shortfall. For the term model, the interval [0.51, 0.69] lies just
below 0.70: on new records its threshold flags somewhat more non-ads than
intended.
